\answer{ex:18:1}{(a)~$0.9949\le w\le1.0049$、すなわち$|1-w|\lesssim0.005$。(b)~$K\ge400$個のシナプス。}
\answer{ex:18:2}{$(\omega_R\tau_w)^2=\sqrt{\alpha(\alpha+2+2\varrho)}/\varrho-1$、$\varrho=\tau_m/\tau_w$。$f_R\approx11.1$~Hz、$|Z(\omega_R)|/|Z(0)|=4.6$。}
\answer{ex:18:3}{(a)~$f=1/\bigl(\tau\ln\frac{\mu}{\mu-\theta}\bigr)$。$1$~Hzには$\mu/\theta-1\approx3.7\cdot10^{-44}$が必要。(b)~$T=\pi\tau/\sqrt\mu$、$f\approx3.2$~Hz。$f\propto\sqrt\mu$である。}
\answer{ex:18:4}{積分器型は$f\lesssim17.7$~Hzで発火し（ローパスフィルタ）、共振器型は$5.5$--$22$~Hzの帯域でだけ発火する（バンドパスフィルタ）。}
\answer{ex:18:5}{(a)~$T_{\min}=\frac{\tau}{w-1}\ln\frac{0.5}{0.3}\approx10.2\ms$。(b)~$B>w-1-0.2=0.3$。}
\answer{ex:18:6}{調整の時定数は$\tau/(\eta\langle r^2\rangle)$、$\eta=\tau\ln10/60\,\text{s}\approx3.8\cdot10^{-3}$。$I\neq0$では重みがずれるので、$e$から指令のコピー$I/\tau$を引く必要がある。}
\answer{ex:18:7}{自由課題。共振器型の有利さは$11$~Hzでわずか$1.35$倍で、$1$~Hzでは$17$分の1と不利になる。組み替えの主な利点は、閾値による帯域の選択にある（問題~\ref{ex:18:4}）。}
